\documentclass[10pt,a4paper]{amsart}
\usepackage[margin=19mm]{geometry}
\usepackage{amsmath,amssymb,mathtools}
\usepackage[scr=rsfs,cal=euler]{mathalfa}
\usepackage{xfrac}
\usepackage[unicode,hidelinks,bookmarks=false]{hyperref}
\newcommand{\ie}{i.e.\ }
\newcommand{\cf}{cf.\ }
\newcommand{\resp}{respectively\ }
\newcommand{\Subsection}{\subsection}
\providecommand{\nobreakdash}{\nobreak\hbox{-}}
\setlength{\emergencystretch}{2em}
\multlinegap0pt
\usepackage{amssymb}
\usepackage{enumerate}
\usepackage{tikz}
\usepackage{float}
\usepackage{mathtools}
\newtheorem{lemma}{Lemma}
\newtheorem{theorem}{Theorem}
\newtheorem{hyp}{Hypothesis}
\DeclareMathOperator{\Irr}{Irr}
\newcommand{\SLi}[3]{\Sigma_{#1,#2}^{#3L}} %math mode required
\DeclareMathOperator{\cd}{cd}
\title{On prime character degree graphs occurring within a family of graphs (iii)}
\author{Mark W. Bissler, Thatcher Debowski, Theodore F. Hoelker, Jacob Laubacher, Lorenzo Ravaglia, G. Sivanesan}
\date{Source: arXiv:2606.05331v1 (CC BY 4.0)}
\begin{document}
\maketitle

\noindent\emph{Source and licence.} On prime character degree graphs occurring within a family of graphs (iii). Version arXiv:2606.05331v1 (CC BY 4.0), distributed by arXiv under the Creative Commons Attribution 4.0 International licence, http://creativecommons.org/licenses/by/4.0/, as recorded in the arXiv licence metadata for this exact version and shown on the abstract page https://arxiv.org/abs/2606.05331v1. Full licence notice: \texttt{licenses/arxiv-2606.05331-CC-BY-4.0.txt}.\par\smallskip

\noindent\textit{Editorial scope.} Counted unit: the article's lemma lem21mLfit (source0.tex lines 632--634). The article's complete original proof of it is delivered with it (source0.tex lines 635--681, one \texttt{proof} environment of 47 lines). No mathematics is written, repaired or re-derived: the statement and the proof are the article's own lines, in the article's own notation, and no retained line is substituted or re-flowed.  Retained uncounted as the source-local context the proof needs: lines 185--187; lines 218--220; lines 493--495; lines 585--587.  Nothing was omitted from any retained span, so the package carries no omission record.  The retained lines cite 7 works, whose entries are transcribed verbatim from the article's own bibliography source in the e-print and delivered with short editorial labels recorded in the item's reference mapping.  No retained line contains a figure environment, an image-inclusion command or drawing code, so no image is delivered and the package declares no asset dependency.  Editorial formatting only, in addition: an AMS \texttt{amsart} preamble that restates the article's own declarations and macros used by the retained lines, namely the environments \texttt{lemma}, \texttt{theorem}, \texttt{hyp} on separate counters, together with the standard packages those lines need.  The retained environments print in this document as Lemma 1, Theorem 1, Hyp 1 in order of appearance, whereas the article prints the same items with the numbering of its own full document.  The complete native source of the article as distributed in the arXiv e-print for this exact version is bundled unchanged as \texttt{sources/202111/source0.tex}.
\begin{lemma}\label{bigoh}\emph{(\cite{BLL2})}
Let $G$ be a solvable group, and suppose $p$ is an admissible vertex of $\Delta(G)$. For every proper normal subgroup $N$ of $G$, suppose that $\Delta(N)$ is a proper subgraph of $\Delta(G)$. Then $O^p(G)=G$.
\end{lemma}

\begin{theorem}\label{KT}\emph{(\cite{BL})}
The graph $\Gamma_{k,t}$ occurs as the prime character degree graph of a solvable group precisely when $t=1$ or $k=t=2$. Otherwise $\Gamma_{k,t}$ does not occur as the prime character degree graph of any solvable group, nor does any proper connected subgraph of $\Gamma_{k,t}$ with the same vertex set whenever $k\geq t\geq2$.
\end{theorem}

\begin{hyp}\label{hyp21mL}
Given any integer $t\geq2$, we assume that the graph $\SLi{2}{1}{t}$ does not occur as the prime character degree graph of any solvable group.
\end{hyp}

\begin{lemma}\label{lem21ad}
Suppose that we are under the conditions of Hypothesis \ref{hyp21mL}. Further suppose that $\SLi{2}{1}{(t+1)}=\Delta(G)$ for some finite solvable group $G$, where $|G|$ is minimal. Then $G$ does not have any normal nonabelian Sylow $p$-subgroup for any $p\in\{a_1,a_2\}\cup\{b_1,b_3\}\cup\{c_1,c_2,\ldots,c_{t+1}\}$. In particular, $p$ is a strongly admissible vertex.
\end{lemma}

\begin{lemma}\label{lem21mLfit}
Suppose that we are under the conditions of Hypothesis \ref{hyp21mL}. Further suppose that $\SLi{2}{1}{(t+1)}=\Delta(G)$ for some finite solvable group $G$, where $|G|$ is minimal. Then the Fitting subgroup $\textit{F}$ of $G$ is minimal normal in $G$.
\end{lemma}

\begin{proof}
Following the aforementioned discussion, $F$ is the Fitting subgroup of $G$ and $H$ is defined as above. Furthermore, we set $E$ to be the Fitting subgroup of $H$. For the sake of contradiction, suppose that $F$ is not minimal normal in $G$. Consequently, we suppose that there exists some normal subgroup $N$ of $G$ such that $1<N<F$. We will show that no such $N$ can exist.

Notice that Theorem III 4.5 of \cite{H} guarantees the existence of another normal subgroup $M$ of $G$ such that $F=N\times M$. Since we know $N<F$ and $F=N\times M$, this then means that $M>1$. These subgroups $N$ and $M$ are now both nontrivial, and so $\rho(G/N)\subset\rho(G)$ and $\rho(G/M)\subset\rho(G)$, where we note that those are proper subsets. In particular, since we now have that $\rho(G)\setminus\rho(G/N)$ is nonempty, then for any prime $q\in\rho(G)\setminus\rho(G/N)$, the group $G/N$ will have a normal nonabelian Sylow $q$-subgroup whose class is a formation. This forces $q\in\rho(G/M)$. The fallout from this is that the vertex set corresponding to the prime character degree graph for $G$ must be the union of the vertex sets for those corresponding to $G/N$ and $G/M$. To wit, $\rho(G)=\rho(G/N)\cup\rho(G/M)$. Following the argument outlined in \cite{LMS}, for instance, we can then conclude that $\rho(G)\setminus[\rho(G/N)\cap\rho(G/M)]$ must lie in a complete subgraph of $\Delta(G)$, and so $\rho(G)\setminus[\rho(G/N)\cap\rho(G/M)]$ must lie in one of the following sets: $A=\{b_1,b_2,b_3\}$, $B=\{a_1,b_1,b_3\}$, $C=\{a_2,b_2\}$, or $D=\{a_1,a_2,c_1,c_2,\ldots,c_{t+1}\}$.

First suppose that $\rho(G)\setminus[\rho(G/N)\cap\rho(G/M)]\subseteq A$. One can follow the argument given in \cite{L3}, for example, and conclude that $E$ is in fact the Hall $A$-subgroup of $H$, and so we can find a character $\chi\in\Irr(G)$ such that $b_1b_2b_3\mid\chi(1)$. Denoting $\theta$ as an irreducible constituent of $\chi_{FE}$, one can verify that $\chi(1)/\theta(1)$ divides $|G:FE|$ and that $\chi(1)$ is relatively prime to $|G:FE|$, which yields $\chi_{FE}=\theta$. Since $b_1$, $b_2$, and $b_3$ all divide $\theta(1)$, and the only possible primes that could divide a character in $\cd(G/FE)$ are the vertices in $\rho(G)\setminus A=\{a_1,a_2,c_1,c_2,\ldots,c_{t+1}\}$, then we get that $\cd(G/FE)$ is trivial due to Gallagher's Theorem, and so $G/FE$ is abelian. This will then force $O^{a_1}(G)<G$, for instance, which is a contradiction, since we know that $a_1$ is admissible by Lemma \ref{lem21ad}, in which case $O^{a_1}(G)=G$ by Lemma \ref{bigoh}. Hence, this case cannot occur.

Next we suppose the case where $\rho(G)\setminus[\rho(G/N)\cap\rho(G/M)]\subseteq B$, which will then yield
$$
\{a_2,b_2,c_1,c_2,\ldots,c_{t+1}\}\subseteq\rho(G/N)\cap\rho(G/M).
$$
Since $b_1\in\rho(G)=\rho(G/N)\cup\rho(G/M)$, then we can, without loss of generality, suppose that $b_1\in\rho(G/N)$. Since we also know that $\rho(G/N)$ is proper in $\rho(G)$, we then have the following three options for the vertex set $\rho(G/N)$: (i) $\rho(G/N)=\{a_2,b_1,b_2,c_1,c_2,\ldots,c_{t+1}\}$, (ii) $\rho(G/N)=\{a_1,a_2,b_1,b_2,c_1,c_2,\ldots,c_{t+1}\}$, or (iii) $\rho(G/N)=\{a_2,b_1,b_2,b_3,c_1,c_2,\ldots,c_{t+1}\}$.

Under case (i), one can see that no connected subgraph with this vertex could occur, because it has two cut vertices (see \cite{MQ}). Therefore, the only graph that could occur is the disconnected subgraph with complete components on the vertex sets $\{a_2,c_1,c_2,\ldots,c_{t+1}\}$ and $\{b_1,b_2\}$. Following \cite{L}, we then know that $G/N$ has a central Sylow $a_1$-subgroup or a central Sylow $b_3$-subgroup, which implies either $O^{a_1}(G)<G$ or $O^{b_3}(G)<G$. This is a contradiction since both $a_1$ and $b_3$ were shown to be admissible, and so $O^{a_1}(G)=G$ and $O^{b_3}(G)=G$, again by Lemma \ref{bigoh}.

Next, suppose (ii), in which case we either end up with a connected proper subgraph of $\Gamma_{t+3,2}$, which does not occur by the addendum of Theorem \ref{KT}, or we have a disconnected subgraph resulting in a central Sylow $b_3$ subgroup of $G/N$, which cannot happen as $O^{b_3}(G)=G$ as above.

Finally for (iii), it is not hard to see that the only subgraph that could occur is disconnected (whereby the graphs that are connected would have two cut vertices, which cannot happen again by \cite{MQ}), resulting in $G/N$ having a central Sylow $a_1$-subgroup, which cannot happen as $a_1$ is admissible. Thus, this case also cannot occur.

We next consider the case where we suppose $\rho(G)\setminus[\rho(G/N)\cap\rho(G/M)]\subseteq C$. Notice that this implies that
$$
\{a_1,b_1,b_3,c_1,c_2,\ldots,c_{t+1}\}\subseteq\rho(G/N)\cap\rho(G/M).
$$
Since $\rho(G/N)$ and $\rho(G/M)$ are both proper in $G$, and because $\rho(G)=\rho(G/N)\cup\rho(G/M)$, then without loss of generality, we must have that
$$
\rho(G/N)=\{a_1,a_2,b_1,b_3,c_1,c_2,\ldots,c_{t+1}\}
$$
and
$$
\rho(G/M)=\{a_1,b_1,b_2,b_3,c_1,c_2,\ldots,c_{t+1}\}.
$$
However, one can check that no connected subgraph with the vertex set $\rho(G/M)$ occurs, and therefore the disconnected subgraph with complete components on the vertex sets $\{a_1,c_1,c_2,\ldots,c_{t+1}\}$ and $\{b_1,b_2,b_3\}$ must occur. Again following the results from \cite{L}, this then yields a central Sylow $a_2$-subgroup of $G/M$, which implies $O^{a_2}(G)<G$. This cannot happen, as $a_2$ is admissible by way of Lemma \ref{lem21ad}. Hence, this case also cannot happen.

Finally, we must therefore suppose that $\rho(G)\setminus[\rho(G/N)\cap\rho(G/M)]\subseteq D$. In this case, notice that this forces $\{b_1,b_2,b_3\}$ to be in both $\rho(G/N)$ and $\rho(G/M)$. Without loss of generality, we may also suppose $c_1\in\rho(G/N)$. Before we look at the cases for what $\rho(G/N)$ can be, we must first address the symmetry of all the vertices $c_i$ for all $1\leq i\leq t+1$. Since these vertices all behave the same, we can arrange them in order, which will help reduce our subsequent cases. We therefore fix the notation of letting $C_i^*:=\{c_2,c_3,\ldots,c_i\}$ for all $2\leq i\leq t+1$. With this in hand, we have eight cases for $\rho(G/N)$: (i) $\{b_1,b_2,b_3,c_1\}$, (ii) $\{a_1,b_1,b_2,b_3,c_1\}$, (iii) $\{a_2,b_1,b_2,b_3,c_1\}$, (iv) $\{a_1,a_2,b_1,b_2,b_3,c_1\}$, (v) $\{b_1,b_2,b_3,c_1\}\cup C_i^*$, (vi) $\{a_1,b_1,b_2,b_3,c_1\}\cup C_i^*$, (vii) $\{a_2,b_1,b_2,b_3,c_1\}\cup C_i^*$, or (viii) $\{a_1,a_2,b_1,b_2,b_3,c_1\}\cup C_i^*$ for $i\neq t+1$. We note that taking $i=t+1$ in the final case would imply that $\rho(G/N)=\rho(G)$, which we know does not happen. Next, we address every case.

Supposing case (i), we have that $\rho(G/N)=\{b_1,b_2,b_3,c_1\}$, and the only graph that arises is the disconnected subgraph with complete components with vertex sets $\{b_1,b_2,b_3\}$ and $\{c_1\}$. Again using \cite{L}, we then have that $G/N$ has a central Sylow $p$-subgroup for some $p\in\{a_1,a_2,c_2,\ldots,c_{t+1}\}$, in which case $O^p(G)<G$. This is our contradiction, as Lemma \ref{bigoh} forces $O^p(G)=G$ for all such relevant $p$ because they were shown to be admissible in Lemma \ref{lem21ad}. Therefore, this case cannot occur.

Next, for (ii), the only graph that arises with this vertex set is disconnected with complete components with vertex sets $\{a_1,c_1\}$ and $\{b_1,b_2,b_3\}$. We get the same contradiction as above.

For both cases (iii) and (iv), no graph (connected or disconnected) occurs with these vertex sets.

Supposing case (v), we see that no such connected graph occurs, but the disconnected graph with complete components with vertex sets $\{b_1,b_2,b_3\}$ and $\{c_1\}\cup C_i^*$ do arise. This, however, is handled identically to how case (i) transpired. This is mainly due to the fact that every vertex outside of $\rho(G/N)$ is admissible.

Case (vi) falls just like case (ii). Likewise, case (vii) behaves like case (iii), and case (viii) is tamed just like case (iv).

Hence, none of these cases can occur, which means our original supposition must have been false, meaning that no such $N$ can occur. Thus, the Fitting subgroup $F$ must be minimal normal in $G$.
\end{proof}

\begin{thebibliography}{99}
\bibitem[BL]{BL} Mark W. Bissler and Jacob Laubacher. Classifying families of character degree graphs of solvable groups. \emph{Int. J. Group Theory}, 8(4):37--46, 2019.
\bibitem[BLL2]{BLL2} Mark W. Bissler, Jacob Laubacher, and Mark L. Lewis. A family of graphs that cannot occur as character degree graphs of solvable groups. \emph{Beitr. Algebra Geom.}, 66(3):727--738, 2025.
\bibitem[H]{H} B. Huppert. \emph{Endliche gruppen. I}. Die Grundlehren der Mathematischen Wissenschaften, Band 134. Springer-Verlag, Berlin-New York, 1967.
\bibitem[L]{L} Mark L. Lewis. Solvable groups whose degree graphs have two connected components. \emph{J. Group Theory}, 4(3):255--275, 2001.
\bibitem[L3]{L3} Mark L. Lewis. Classifying character degree graphs with 5 vertices. In \emph{Finite groups 2003}, pages 247--265. Walter de Gruyter, Berlin, 2004.
\bibitem[LMS]{LMS} Jacob Laubacher, Mark Medwid, and Dylan Schuster. Classifying character degree graphs with seven vertices. \emph{Adv. Group Theory Appl.}, 22:79--121, 2025.
\bibitem[MQ]{MQ} Mark L. Lewis and Qingyun Meng. Solvable groups whose prime divisor character degree graphs are 1-connected. \emph{Monatsh. Math.}, 190(3):541--548, 2019.
\end{thebibliography}
\end{document}
